\section{Divisibility, primes, and finite arithmetic}
\label{sec:elementary}

This chapter develops the integer arithmetic needed by the later chapters.
Every search is bounded.  A statement about an arbitrary primitive recursive
function is a schema with that function fixed, or, equivalently for the
finite assertions used here, a statement about a coded finite table.  It
never means that arithmetic has a total interpreter for all primitive
recursive programs.  The conventional topics correspond to the elementary
portion of the source article and to the university course outline
\cite{mit781}; the arguments and arithmetic bounds below are supplied in full.

We use natural numbers including zero.  A signed integer is represented by
its sign and absolute value; addition, multiplication, comparison, and
absolute value have their usual primitive recursive definitions on these
codes.  For natural numbers, $d\mid n$ means that $n=dq$ for some $q\le n$,
with the separate convention $0\mid0$.  For integers one uses a signed
quotient of absolute value at most $|n|$ when $d\ne0$.  Thus the divisibility
predicate is decidable by a bounded test.  A bounded maximum, minimum,
conjunction, sum, or product is implemented by a fold over its displayed
finite range.  For example, primality is the bounded predicate
\[
 \operatorname{Prime}(p)\ \Longleftrightarrow\
 p\ge2\ \text{ and }\ (\forall d\le p)
       (d\mid p\Longrightarrow d=1\text{ or }d=p).
\]
A bounded universal quantifier in a loop invariant abbreviates a Boolean
fold.  Induction on that Boolean assertion is quantifier-free induction.

\subsection{Division and the Euclidean algorithm}

\begin{theorem}[Division with a bounded algorithm]\label{el:division}
For $a\in\N$ and $b>0$ there are unique natural numbers $q,r$ such that
$a=bq+r$ and $r<b$.  They satisfy $q\le a$ and $r\le a$, and are primitive
recursive functions of $a,b$.
\end{theorem}
\begin{proof}
Start with $(q_0,r_0)=(0,0)$.  To pass from $i$ to $i+1$, replace
$(q_i,r_i)$ by $(q_i,r_i+1)$ if $r_i+1<b$, and by $(q_i+1,0)$ otherwise.
The invariant is
\[
 i=bq_i+r_i,\qquad r_i<b,\qquad q_i\le i,\quad r_i\le i.
\]
In the second branch $r_i<b$ and the failure of $r_i+1<b$ give
$r_i+1=b$.  This proves preservation by direct arithmetic.  After exactly
$a$ steps the required pair has been obtained.  For uniqueness, if
$bq+r=bq'+r'$ and $q<q'$, then
$bq+r<b(q+1)\le bq'\le bq'+r'$, a contradiction.  Decidable trichotomy
therefore gives $q=q'$ and then $r=r'$.  No least-number principle is used.
\end{proof}

We write these outputs as $\lfloor a/b\rfloor$ and $a\bmod b$.
For a signed integer $z$ the same construction on $|z|$, followed when
necessary by $(-q-1,b-r)$, gives a unique residue in $[0,b)$.

\begin{theorem}[Euclidean algorithm and B\'ezout identity]
\label{el:bezout}
There are primitive recursive functions $g(a,b)\in\N$ and
$u(a,b),v(a,b)\in\Z$ satisfying
\[
 g\mid a,\quad g\mid b,\quad ua+vb=g,
 \qquad d\mid a\text{ and }d\mid b\Longrightarrow d\mid g.
\]
Here $g(0,0)=0$; otherwise $g>0$.  Consequently $g$ is the greatest
common divisor.  A computation with at most $a+b+1$ Euclidean steps
suffices, and the displayed coefficients may be chosen with
\[
 |u|,|v|\le (a+b+2)^{a+b+1}.
\]
\end{theorem}
\begin{proof}
Maintain two remainders $x,y$ and two coefficient rows $U,V$, initially
$x=a$, $y=b$, $U=(1,0)$, $V=(0,1)$.  If $y=0$, keep the state fixed.
Otherwise put $q=\lfloor x/y\rfloor$ and replace
\[
 (x,y,U,V)\quad\text{by}\quad(y,x-qy,V,U-qV).
\]
The two remainders remain the linear combinations described by the rows.
For any fixed $d$, the conditions $d\mid x,d\mid y$ are equivalent to
$d\mid y,d\mid x-qy$: one direction is subtraction and the other is
$x=(x-qy)+qy$.  In particular the initial common divisors and the current
common divisors coincide.

Whenever a nonzero second remainder is replaced, its successor is a
smaller natural number.  After the first step both remainders are at most
$\max(a,b)$; the second remainder starts at $b$ and decreases by at least
one on every active step.  A loop with $a+b+1$ steps has therefore reached
$y=0$.  A quantifier-free termination invariant is: either $y_i=0$, or
$y_i+i\le b$ .
At the terminal state let $g=x$ and let $(u,v)=U$.  The common-divisor
invariant gives all the claimed divisibilities, including $g\mid a,b$.
If the initial pair is nonzero then an active step never makes both
remainders zero, so $g>0$.  A positive common divisor of $a,b$ divides
the positive number $g$ and is therefore at most $g$.

Every quotient is at most $a+b$.  If all four coefficient magnitudes are
at most $C$, a step makes them at most $(a+b+1)C$.  Starting with $C=1$
and using the larger base $a+b+2$ gives the asserted bound.  All invariants
are equalities, inequalities, and bounded divisibility tests.  One may
preserve the common-divisor assertion for each fixed $d$ separately;
there is no need to induct on an unbounded universal formula.
\end{proof}

\begin{theorem}[Linear Diophantine equations]\label{el:linear}
Let $a,b\in\N$, not both zero, and $c\in\Z$, and put $g=g(a,b)$.
The equation $ax+by=c$ has integer solutions exactly when $g\mid c$.
When it does, one solution is
\[
 x_0=u(a,b)(c/g),\qquad y_0=v(a,b)(c/g).
\]
All solutions are
\[
 x=x_0+(b/g)t,\qquad y=y_0-(a/g)t\qquad(t\in\Z).
\]
For $a=b=0$ solvability is exactly $c=0$.
\end{theorem}
\begin{proof}
A common divisor divides every linear combination, proving necessity.
Multiplying the B\'ezout identity by $c/g$ proves sufficiency.  Divide
$a,b$ by $g$ and write $a',b'$; their gcd is one, since a common divisor
$h>1$ would make $gh$ a common divisor of $a,b$ greater than $g$.
If $(x,y)$ is a solution, then
$a'(x-x_0)=-b'(y-y_0)$.  For coprime integers $A,B$, a B\'ezout
identity $rA+sB=1$ shows that $B\mid Az$ implies $B\mid z$:
multiply the identity by $z$.  If $b'>0$, this proves
$x-x_0=b't$, and cancellation gives $y-y_0=-a't$.
If $b'=0$, then $a'=1$, so $x=x_0$ and $y$ is arbitrary; the same
formula still describes every solution.  The constructed solution has
coordinate magnitudes at most
$(a+b+2)^{a+b+1}|c|$.  The assertion about an arbitrary proposed solution
uses its own coordinates as parameters, not an unbounded search for $t$.
\end{proof}

\subsection{Prime decomposition and its consequences}

\begin{lemma}[A least divisor is prime]\label{el:leastprime}
Every $n\ge2$ has a prime divisor at most $n$, obtained by taking the
least $d\in[2,n]$ dividing $n$.  If $n$ is composite, it has a prime
divisor $p$ with $p^2\le n$.
\end{lemma}
\begin{proof}
The search succeeds because $n\mid n$.  If its output $d$ were
composite, it would have a divisor $e$ with $1<e<d$; then $e\mid n$,
contrary to the checked minimality.  If $n$ is composite and $p$ is its
least prime divisor, write $n=pk$ with $k>1$.  A prime divisor of $k$
is at least $p$, so $k\ge p$ and $p^2\le n$.
All alternatives are decidable within $[0,n]$.
\end{proof}

\begin{theorem}[Euclid's lemma]\label{el:euclidlemma}
For a prime $p$, if $p\mid ab$, then $p\mid a$ or $p\mid b$.
Consequently, if $\gcd(a,m)=1$ and $m\mid ab$, then $m\mid b$.
\end{theorem}
\begin{proof}
Test $p\mid a$.  If it fails, the positive divisor $g(p,a)$ of $p$
cannot equal $p$, so it equals one.  A B\'ezout identity $up+va=1$,
multiplied by $b$, expresses $b$ as a sum of two multiples of $p$.
The coprime version follows by precisely the same multiplication.
This proof uses excluded middle only for the bounded predicate
$p\mid a$, whose characteristic function is already primitive recursive.
\end{proof}

\begin{theorem}[Fundamental theorem of arithmetic]
\label{el:fta}
Every $n\ge1$ has a primitive recursively computed nondecreasing list
of primes $p_0,\ldots,p_{\ell-1}$, with $\ell\le n$, whose product is $n$.
Any other nondecreasing prime list with product $n$ is identical to it.
Equivalently, the multiset of prime factors is unique.
\end{theorem}
\begin{proof}
Keep a remainder $r$, initially $n$, and an initially empty factor list.
If $r>1$, append its least prime divisor and divide $r$ by that divisor;
if $r=1$, do nothing.  An active step strictly decreases $r$, and the
invariant is
\[
 n=r\prod_{j<\ell}p_j,\qquad r\ge1,\qquad
 \ell\le i,\qquad (r>1\Longrightarrow r+i\le n).
\]
A loop of $n$ steps therefore ends with $r=1$.  A prime factor of the
new remainder was also a divisor of the old remainder, so it is at least
the prime just removed; the list is nondecreasing.  Every list entry is
at most $n$, and a length-$n$ array in base $n+1$, together with its
length, is a primitive recursive code.

For uniqueness, the first prime of one nonempty list divides the
product of the other.  Repeated use of Euclid's lemma, over at most the
length of that list, gives an entry divisible by this prime.  Since that
entry is prime, the entries are equal.  Interchanging the lists shows
that their first entries are equal: each is at most an entry equal to
the first entry of the other list.  Cancel that positive common entry
and repeat.  A prime product is at least $2^{\ell}$, so a list with
product one must be empty.  Thus neither list can end earlier than the
other.  This is an induction on a bounded comparison loop, with the
invariant that the canceled prefixes agree and the remaining products
agree.  It requires no induction on an existential factorization claim.
\end{proof}

For $p\ge2$ and $n\ge1$ define
\[
 v_p(n)=\max\{e\le n:p^e\mid n\}.
\]
The bound is valid because $p^e\ge e+1$ by induction.  This definition is
used only with $p$ prime when the valuation laws are invoked.

\begin{theorem}[Valuations, gcd, and lcm]\label{el:valuation}
For primes $p$ and positive $a,b$,
\[
 v_p(ab)=v_p(a)+v_p(b),\qquad
 a\mid b\ \Longleftrightarrow\
       (\forall p\le\max(a,b))
       (\operatorname{Prime}(p)\Longrightarrow v_p(a)\le v_p(b)).
\]
The valuations of $\gcd(a,b)$ and $\operatorname{lcm}(a,b)$ are the
respective minima and maxima.  In particular
\[
 \gcd(a,b)\operatorname{lcm}(a,b)=ab.
\]
\end{theorem}
\begin{proof}
Collect the occurrences of $p$ in the unique prime lists.  Concatenating
and sorting the lists for $a$ and $b$ gives a prime list for $ab$, so
uniqueness says that the occurrences add.  If $a\mid b$, write $b=ac$
and use additivity.  Conversely, subtract the exponent of each prime
in $a$ from that in $b$ and multiply the resulting nonnegative prime
powers.  The product $c$ satisfies $ac=b$, proving divisibility.
Take the products with minimum and maximum exponents to construct $G,L$.
The exponent criterion proves that $G$ divides both inputs and every
common divisor divides $G$; similarly $L$ is a common multiple dividing
every common multiple.  Thus $G$ is the Euclidean gcd and $L$ is the
least positive common multiple.  Since
$\min(r,s)+\max(r,s)=r+s$, their product is $ab$.
All prime scans can be limited to $\max(a,b)$; any proposed common
multiple is handled with its own finite bound.  The output $L\le ab$.
\end{proof}

Integer fractions use the primitive recursive canonical reduction proved in
Theorem~\ref{dio:normal}.  The normalization map sends $(a,b)$, $b>0$,
to $(a/g,b/g)$ with $g=\gcd(|a|,b)$; all fraction calculations in this
chapter may use these canonical finite codes.

\begin{theorem}[Euclid's theorem with a bound]\label{el:euclid}
For every $N\in\N$ a prime $p$ with
\[
             N<p\le N!+1
\]
can be computed primitive recursively.  Consequently there is no finite
list containing all primes.
\end{theorem}
\begin{proof}
Take the least prime divisor of $N!+1\ge2$.  If this prime were at most
$N$, it would divide both $N!$ and $N!+1$, and hence divide one, a
contradiction.  For a supplied finite prime list, use $N$ equal to its
maximum (zero for the empty list).  The factorial is defined by primitive
recursion, so this is an explicit uniform witness rather than an
unbounded search for the next prime.
\end{proof}

\begin{theorem}[Correctness of the sieve of Eratosthenes]
\label{el:sieve}
There is a primitive recursive Boolean array marking exactly the primes
at most $N$.  The usual sieve, which crosses off multiples of $d$ from
$d^2$ onward whenever $d$ is still unmarked, is correct.
\end{theorem}
\begin{proof}
Initially mark all entries $2,\ldots,N$ as candidates.  Scan
$d=2,\ldots,\lfloor\sqrt N\rfloor$, where the square root is the bounded
maximum of $r\le N$ with $r^2\le N$.  If $d$ is a candidate, cross off
$d^2+kd\le N$ for $k\le N$.
Inductively a candidate $d$ must be prime: if it were composite, its
least prime divisor $q<d$ would have crossed it off, because $d\ge q^2$.
No prime is ever crossed off, since every crossed-off integer is a
product of two integers at least two.  Finally, every composite $n\le N$
has a least prime divisor $p$ with $p^2\le n$; the step $d=p$ occurs,
$p$ is still a candidate, and $n$ is crossed off.  The array can be coded
by a number below $2^{N+1}$.  At most $N^2$ crossing-off checks and $N$
outer steps are enough.  The invariant ``a crossed-off entry has a known
proper factor, and every processed prime has crossed off its prescribed
multiples'' is a bounded Boolean assertion on that array.
\end{proof}

\subsection{Congruences, finite residue systems, and powers}

For $m>0$ put $a\equiv b\pmod m$ when $a\bmod m=b\bmod m$.
Division identifies this with $m\mid(a-b)$.  Substitution in differences
proves reflexivity, symmetry, transitivity, and compatibility with
addition and multiplication.  B\'ezout's identity gives the useful rule
\[
 \gcd(a,m)=1,\quad ax\equiv ay\pmod m
       \quad\Longrightarrow\quad x\equiv y\pmod m.
\]
All residue representatives below lie in a stated finite interval.

\begin{theorem}[Solving and counting linear congruences]
\label{el:linearcong}
Let $m>0$, $a\in\N$, $b\in\Z$, and $g=\gcd(a,m)$.
The congruence $ax\equiv b\pmod m$ is solvable exactly when $g\mid b$.
When solvable, it has exactly $g$ representatives in $[0,m)$, namely
\[
 x_0+k(m/g)\qquad(0\le k<g),
\]
where $0\le x_0<m/g$ is computed by a B\'ezout inverse modulo $m/g$.
\end{theorem}
\begin{proof}
A solution gives $ax-my=b$, proving necessity.  If $g\mid b$, divide
by $g$ to obtain $a'x\equiv b'\pmod {m'}$ with
$\gcd(a',m')=1$.  Compute $u$ with $ua'+vm'=1$ and take
$x_0=(ub')\bmod m'$.  Multiplication by $u$ shows that every solution
is congruent to $x_0$ modulo $m'$.  The integers in $[0,m)$ with that
residue are exactly $x_0+km'$ for $k<g$, by division of a proposed
representative by $m'$.  They are distinct because their differences
have magnitude less than $m$ and a nonzero difference of their indices
cannot be a multiple of $g$.  This includes $m'=1$ and $a=0$.
\end{proof}

\begin{theorem}[Generalized Chinese remainder theorem for two moduli]
\label{el:crt2}
For $m,n>0$, the system
$x\equiv a\pmod m$, $x\equiv b\pmod n$ has a solution exactly when
$\gcd(m,n)\mid(b-a)$.  There is then a unique solution in
$[0,L)$, where $L=\operatorname{lcm}(m,n)\le mn$.
\end{theorem}
\begin{proof}
Necessity follows by subtracting the congruences.  Put $g=\gcd(m,n)$,
$m'=m/g$, $n'=n/g$.  If $g\mid b-a$, solve
$m't\equiv(b-a)/g\pmod {n'}$ using the preceding theorem.  The integer
$a+mt$ satisfies both congruences.  Reduce it modulo $L=mn'$ to obtain
the required representative; both moduli divide $L$, so reduction
preserves the congruences.  Two solutions differ by a multiple of both
$m$ and $n$, and hence by a multiple of $L$ by the valuation criterion.
Two representatives in $[0,L)$ with that property are equal.
\end{proof}

\begin{theorem}[Chinese remainder theorem for a finite system]
\label{el:crt}
For positive moduli $m_0,\ldots,m_{k-1}$ and integer residues $a_i$,
a simultaneous solution exists exactly when
\[
 \gcd(m_i,m_j)\mid(a_i-a_j)\qquad(i,j<k).
\]
When it exists, there is a unique representative below
$L=\operatorname{lcm}(m_0,\ldots,m_{k-1})$, and it is primitive
recursively computable.  In particular pairwise coprime moduli admit
arbitrary residues, with $L=\prod_i m_i$.
\end{theorem}
\begin{proof}
First prove the pairwise coprime case by successively applying
Theorem~\ref{el:crt2}.  The product of the previously used moduli is
coprime to the next: otherwise a prime dividing their gcd would divide
one previous modulus, contradicting the hypothesis.  Thus each step
is possible and gives one representative below the enlarged product.

For general moduli take $M=\prod_i m_i$ (the empty product is one).
For each prime $p\le M$ occurring in some modulus, set
$e_p=\max_{i<k}v_p(m_i)$ and select the least index $j_p$ attaining it.
The compatibility hypotheses imply
\[
 a_{j_p}\equiv a_i\pmod {p^{v_p(m_i)}}
\]
for every $i$: the indicated prime power divides
$\gcd(m_i,m_{j_p})$.  Use the coprime case to solve
$x\equiv a_{j_p}\pmod {p^{e_p}}$ for the finitely many occurring primes.
Its modulus is $L=\prod_p p^{e_p}\le M$.  For each $i$, the difference
$x-a_i$ is divisible by every prime power in $m_i$, hence by $m_i$.
Necessity follows by subtraction, and uniqueness follows because a
difference divisible by all moduli is divisible by their lcm.
All maxima, index choices, and prime scans are bounded by $k$ and $M$.
For $k=0$ the assertion has the unique representative zero modulo one.
\end{proof}

Define
\[
 \varphi(n)=\#\{1\le r\le n:\gcd(r,n)=1\}\qquad(n\ge1).
\]
For $n>1$ the same count uses the residues $1,\ldots,n-1$.
The inverse of a coprime residue is computed by B\'ezout or by a scan
of $[0,n)$, so inverse and multiplication are primitive recursive maps
on an explicitly finite list.

\begin{theorem}[Euler's theorem]\label{el:euler}
If $n\ge1$ and $\gcd(|a|,n)=1$, then
\[
                    a^{\varphi(n)}\equiv1\pmod n.
\]
\end{theorem}
\begin{proof}
For $n=1$ every pair of integers is congruent.  For $n>1$, enumerate
the coprime residues as $r_0<\cdots<r_{h-1}$, where $h=\varphi(n)$.
Multiplication by $a$, followed by reduction, sends this list to itself:
a common divisor of $ar_i$ and $n$ would, using an inverse for $a$,
divide $r_i$.  An inverse for $a$ also gives the inverse map, so this
is a permutation of the finite list.  Products of a finite list are
unchanged by such a permutation.  Concretely, scan the image list,
find each next entry in the remaining source list, and swap it into
place; every search has length at most $h$ and each swap preserves
the product.  Hence
\[
 a^h\prod_{i<h}r_i\equiv\prod_{i<h}r_i\pmod n.
\]
The product is coprime to $n$: multiply the inverses of the entries.
Coprime cancellation proves the assertion.  The finite permutation
argument needs only bounded arrays and quantifier-free induction.
\end{proof}

\begin{corollary}[Fermat's little theorem]\label{el:fermat}
For a prime $p$, $p\nmid a$ implies $a^{p-1}\equiv1\pmod p$.
For every integer $a$ one has $a^p\equiv a\pmod p$.
\end{corollary}
\begin{proof}
The nonzero residues modulo a prime are exactly its coprime residues,
so $\varphi(p)=p-1$.  Apply Euler's theorem and multiply by $a$.
If $p\mid a$, the second congruence holds directly.  The case split
is the bounded divisibility test.
\end{proof}

\begin{theorem}[Order of a residue]\label{el:order}
For $n>1$ and $\gcd(|a|,n)=1$, the least positive exponent
$d\le\varphi(n)$ with $a^d\equiv1\pmod n$ exists and is primitive
recursive.  For every $t\in\N$,
\[
 a^t\equiv1\pmod n\ \Longleftrightarrow\ d\mid t,
 \qquad\text{and consequently }d\mid\varphi(n).
\]
\end{theorem}
\begin{proof}
Euler's theorem supplies a witness within the search range.  Write
$t=qd+r$, $0\le r<d$.  Since $a^d\equiv1$, one has
$a^t\equiv a^r$.  If this is one and $r>0$, the search defining $d$
would have found $r$ earlier.  Thus $r=0$.  Conversely, a multiple
of $d$ gives a power of $a^d$ and hence residue one.  Substitute
$t=\varphi(n)$ for the last assertion.
\end{proof}

\begin{theorem}[Wilson's theorem and converse]\label{el:wilson}
For every $n\ge2$,
\[
 \operatorname{Prime}(n)\quad\Longleftrightarrow\quad
 (n-1)!\equiv-1\pmod n.
\]
\end{theorem}
\begin{proof}
Let $p$ be prime.  Every residue $1,\ldots,p-1$ has a unique inverse.
An inverse equals the original residue exactly when
$p\mid(x-1)(x+1)$, which, by Euclid's lemma, gives $x=1$ or $x=p-1$
within the specified interval.  For odd $p$ these are distinct and
all other entries form pairs with product one.  To make the pairing
literal, maintain a Boolean ``used'' array; take the least unused
non-fixed entry, pair it with its computed inverse, and mark both.
The inverse is unused, since if it had already appeared then its
partner would have been the present entry.  After at most $p$ steps
all non-fixed entries have been paired.  The factorial residue is
therefore $1\cdot(p-1)=-1$.  For $p=2$, directly $1!=-1\pmod2$.

Conversely, if $n$ is composite choose a proper divisor $d$ with
$1<d<n$.  It divides $(n-1)!$.  If $n$ divided $(n-1)!+1$, then $d$
would also divide that number and hence divide one, a contradiction.
Because primality is decidable, this exclusion of the composite branch
proves primality constructively.
\end{proof}

\begin{theorem}[RSA correctness, including nonunits]\label{el:rsa}
Let $p,q$ be distinct primes, $N=pq$, and let positive $e,d$ satisfy
$ed\equiv1\pmod{(p-1)(q-1)}$.  For every integer message $m$,
\[
                    (m^e)^d\equiv m\pmod N.
\]
All key arithmetic and message transformations here are primitive
recursive.  This assertion makes no claim about cryptographic security.
\end{theorem}
\begin{proof}
Since $ed\ge1$, write $ed=1+t(p-1)(q-1)$ with $t\ge0$.
Modulo $p$, test whether $p\mid m$.  If so, both $m^{ed}$ and $m$
are zero.  Otherwise Fermat's theorem gives
$m^{ed}=m(m^{p-1})^{t(q-1)}\equiv m$.  The same proof works modulo
$q$.  These distinct prime moduli are coprime, so the difference is
divisible by their product.  The exponentiation identity
$(m^e)^d=m^{ed}$ is a direct induction.  To compute encryption or
decryption one can reduce after every multiplication; a loop of
length equal to the exponent keeps all residues below $N$.
\end{proof}

\begin{theorem}[Base-expansion divisibility rules]\label{el:digits}
If $b\ge2$ and $n=\sum_{i<k}d_i b^i$, $0\le d_i<b$, then
\[
 n\equiv\sum_{i<k}d_i\pmod{b-1},\qquad
 n\equiv\sum_{i<k}(-1)^i d_i\pmod{b+1}.
\]
If $m\mid b^j$, divisibility by $m$ depends only on the last $j$
base-$b$ digits.  Thus the usual decimal rules for $3,9,11$, and
powers dividing $10^j$ follow.
\end{theorem}
\begin{proof}
Modulo $b-1$ every power of $b$ is one; modulo $b+1$ its $i$th power
is $(-1)^i$.  These claims follow by induction on $i$ and can be
substituted into the finite sum.  For the last assertion, split the
sum at $j$; every term of the remaining part contains $b^j$.
The digit list itself is obtained by repeated division by $b$ for at
most $n+1$ steps, padding by zeros after the quotient reaches zero.
The invariant says that the original number is the extracted low
part plus $b^i$ times the current quotient.
\end{proof}

\subsection{Factorials, binomial coefficients, and their valuations}

Define $\binom nk$ by Pascal recursion, setting it to zero when $k>n$.
Finite summation of Pascal's rule proves
$(x+y)^n=\sum_{k=0}^n\binom nkx^ky^{n-k}$ by induction on $n$:
multiply the previous sum by $x+y$, shift the second index, and add
adjacent coefficients.  The same recursion proves
$k!(n-k)!\binom nk=n!$ for $k\le n$.  Indeed, at an interior
index $1\le k\le n$, multiplying Pascal's identity for
$\binom{n+1}{k}$ by $k!(n+1-k)!$ gives the sum
$(n+1-k)n!+kn!=(n+1)!$ by the two induction identities.
The boundary indices are immediate.

\begin{theorem}[Prime binomial divisibility and the Frobenius identity]
\label{el:binomial}
For a prime $p$ and $0<k<p$ one has $p\mid\binom pk$.
Consequently $(x+y)^p\equiv x^p+y^p\pmod p$ for integer $x,y$,
and the coefficientwise identity
$(A+B)^p=A^p+B^p$ holds for finite polynomials modulo $p$.
\end{theorem}
\begin{proof}
The factorial product $k!(p-k)!$ contains no factor divisible by $p$,
by repeated Euclid's lemma; it is coprime to $p$.  The factorial
identity shows that $p$ divides this product times $\binom pk$,
because $p\mid p!$.  Cancel the coprime product.  All the inner
coefficients of the binomial expansion therefore vanish modulo $p$.
The same calculation works in the finite polynomial ring: addition
and multiplication are bounded coefficient sums, and the binomial
expansion uses only commutativity and distributivity.  There is no
appeal to an infinite polynomial object.
\end{proof}

\begin{theorem}[Legendre's factorial valuation formula]
\label{el:legendre}
For a prime $p$ and $n\ge0$,
\[
        v_p(n!)=\sum_{j=1}^{n}\left\lfloor\frac{n}{p^j}\right\rfloor.
\]
If $n=\sum_{i<k}d_i p^i$ is its base-$p$ expansion and
$s_p(n)=\sum_{i<k}d_i$, this is equivalently
\[
                 (p-1)v_p(n!)=n-s_p(n).
\]
\end{theorem}
\begin{proof}
Valuation additivity gives $v_p(n!)=\sum_{a=1}^n v_p(a)$.
For each $a\le n$, the exponent $v_p(a)$ counts exactly the indices
$j\in[1,n]$ with $p^j\mid a$: the valid powers form the initial segment
through $v_p(a)$, and $p^n>n$ for $n>0$.  Interchange the two finite
sums.  The number of multiples of $p^j$ in $[1,n]$ is
$\lfloor n/p^j\rfloor$, by division, proving the first formula.
For the second, division of the base-$p$ expansion gives
$\lfloor n/p^j\rfloor=\sum_{i\ge j}d_i p^{i-j}$.
Summing and using
$(p-1)\sum_{h=0}^{i-1}p^h=p^i-1$ gives
$(p-1)v_p(n!)=\sum_i d_i(p^i-1)=n-s_p(n)$.
For $n=0$ all sums are empty and $0!=1$ has valuation zero.
\end{proof}

\begin{theorem}[Lucas's binomial congruence]\label{el:lucas}
Let $p$ be prime and write $n=\sum_i n_i p^i$ and
$k=\sum_i k_i p^i$ with digits in $[0,p)$, padding with zeros to the
same finite length.  Then
\[
                    \binom nk\equiv\prod_i\binom{n_i}{k_i}\pmod p,
\]
where a binomial coefficient with lower index exceeding the upper
is zero.
\end{theorem}
\begin{proof}
Iterate the finite polynomial identity of
Theorem~\ref{el:binomial} to obtain
$(1+X)^{p^i}\equiv1+X^{p^i}$ coefficientwise modulo $p$.
Exponent addition therefore gives
\[
 (1+X)^n\equiv\prod_i(1+X^{p^i})^{n_i}\pmod p.
\]
Expanding the right side chooses indices $j_i$ with
$0\le j_i\le n_i<p$, contributing the exponent
$\sum_i j_i p^i$ and coefficient $\prod_i\binom{n_i}{j_i}$.
Uniqueness of base-$p$ division says that this exponent is $k$ exactly
when $j_i=k_i$ for every $i$.  If some digit $k_i>n_i$ no term exists
and the coefficient is zero; otherwise precisely one digit tuple
exists, yielding the formula.  This coefficient comparison is finite:
arrays of length $n+1$ suffice, with integer coefficients bounded by
$2^n$, and digit lists have length at most $n+k+1$.  For $k>n$ the
same comparison uses the zero coefficient outside degree $n$.
\end{proof}

\begin{theorem}[Kummer's carry formula]\label{el:kummer}
For a prime $p$ and $0\le k\le n$, the exponent of $p$ in
$\binom nk$ equals the number of carries when adding $k$ and $n-k$
in base $p$.  Equivalently,
\[
 (p-1)v_p\!\left(\binom nk\right)
     =s_p(k)+s_p(n-k)-s_p(n).
\]
\end{theorem}
\begin{proof}
The binomial coefficient is positive for $k\le n$, as follows from
Pascal recursion.  Taking valuations in
$k!(n-k)!\binom nk=n!$ and applying Legendre's formula gives the
displayed digit-sum identity.  To interpret it, let $a_i,b_i,n_i$
be the base-$p$ digits of $k,n-k,n$, padded with zeros, and let
$c_i$ be the carry into position $i$, with $c_0=0$.
The grade-school addition recursion is the division of
$a_i+b_i+c_i$ by $p$.  Since $a_i,b_i<p$ and $c_i\le1$, each carry
$c_{i+1}$ is zero or one, and
\[
                 a_i+b_i+c_i=n_i+p c_{i+1}.
\]
There are at most $n+2$ relevant positions, and the last carry is
zero after padding because the sum is $n$.  Sum these equations
without positional weights.  The sums of incoming and outgoing
carries agree, leaving
$s_p(k)+s_p(n-k)-s_p(n)=(p-1)\sum_i c_{i+1}$.
Cancel $p-1>0$.  The carry array is computed by a fixed bounded loop,
so both the equality and the count use only quantifier-free induction.
\end{proof}

\subsection{Arithmetic functions and finite inversion}

For $n\ge1$ and $r\in\N$ define
\[
 \tau(n)=\sum_{d=1}^n[d\mid n],\qquad
 \sigma_r(n)=\sum_{d=1}^n[d\mid n]d^r,\qquad \sigma=\sigma_1.
\]
Here $[P]$ is one or zero according as the decidable predicate $P$
holds or fails.  Define $\mu(n)=0$ when a square of a prime divides
$n$; otherwise define $\mu(n)=(-1)^t$, where $t$ is the number of
its distinct prime factors.  In particular $\mu(1)=1$.
An arithmetic function is called multiplicative if its value at one
is one and it satisfies $f(ab)=f(a)f(b)$ whenever $a,b$ are positive
and coprime.  Statements involving an unspecified $f$ below are
fixed-function schemas, or statements about supplied finite tables.

\begin{theorem}[Divisor formulas and multiplicativity]
\label{el:divisorfunc}
If $n=\prod_{i<t}p_i^{e_i}$ with distinct primes and $e_i>0$, then
\[
 \tau(n)=\prod_{i<t}(e_i+1),\qquad
 \sigma_r(n)=\prod_{i<t}\left(\sum_{j=0}^{e_i}p_i^{rj}\right).
\]
The functions $\tau$ and $\sigma_r$ are multiplicative.  For $r=1$,
\[
             (p-1)\sigma(p^e)=p^{e+1}-1.
\]
\end{theorem}
\begin{proof}
The valuation criterion gives a bijection between positive divisors
$d$ of $n$ and tuples $(j_0,\ldots,j_{t-1})$ with $0\le j_i\le e_i$:
the divisor is $\prod_i p_i^{j_i}$.  Counting tuples gives the formula
for $\tau$.  Summing $d^r=\prod_i p_i^{rj_i}$ and distributing the
finite product gives the second formula.  For coprime $a,b$ their
prime supports are disjoint, so both products concatenate; this proves
multiplicativity.  The last equation follows by subtracting the finite
geometric sum from $p$ times that sum.  All tuple codes are bounded by
$\prod_i(e_i+1)\le(n+1)^n$; direct divisor scans use only $n$ steps.
\end{proof}

\begin{theorem}[Squarefree decomposition and parity of divisor functions]
\label{el:squareparity}
Every $n\ge1$ has unique positive integers $s,t$ such that
$n=st^2$ and $s$ is squarefree (no prime square divides it).
Both are primitive recursively computable and at most $n$.
Furthermore,
\[
 \tau(n)\text{ is odd}\ \Longleftrightarrow\ n\text{ is a square},
\]
\[
 \sigma(n)\text{ is odd}\ \Longleftrightarrow\
 n\text{ is a square or twice a square}.
\]
\end{theorem}
\begin{proof}
For every prime exponent $e$ in the factorization of $n$, divide it
as $e=2q+r$ with $r\in\{0,1\}$.  Put the factor $p^r$ into $s$ and
$p^q$ into $t$.  Unique factorization proves $n=st^2$, squarefreeness,
and uniqueness of these choices.  The product formula
$\tau(n)=\prod_p(e_p+1)$ is odd exactly when every $e_p$ is even,
which is exactly the condition for $n$ to be a square.
For an odd prime $p$, the sum $1+p+\cdots+p^e$ has parity $e+1$;
for $p=2$ it is always odd.  Thus $\sigma(n)$ is odd exactly when
every odd-prime exponent is even.  In that case the exponent of
two is either even, giving a square, or odd, giving twice a square.
The converse follows from the same parity calculation.
All root witnesses are the displayed products of half-exponents,
so the existential conclusions carry explicit bounded constructions.
\end{proof}

\begin{theorem}[Totient formula and multiplicativity]
\label{el:totient}
For a prime $p$ and $e\ge1$,
$\varphi(p^e)=p^{e-1}(p-1)$.  The totient is multiplicative, and hence
\[
             \varphi(n)=\prod_{p^e\parallel n}p^{e-1}(p-1)
             \qquad(n\ge1).
\]
\end{theorem}
\begin{proof}
A residue modulo $p^e$ is noncoprime exactly when it is divisible by
$p$, using prime factorization.  Among $1,\ldots,p^e$ precisely
$p^{e-1}$ numbers are multiples of $p$.  Subtracting their count proves
the prime-power formula.  If $a,b$ are coprime, the Chinese remainder
map from residues modulo $ab$ to pairs of residues modulo $a,b$ is
bijective.  A residue is coprime to $ab$ exactly when each component
is coprime to its modulus: a prime common divisor would divide $a$
or $b$.  Restricting the finite bijection to these decidable sublists
gives $\varphi(ab)=\varphi(a)\varphi(b)$.  Factorization then proves
the product formula; the empty product gives $\varphi(1)=1$.
\end{proof}

\begin{theorem}[The divisor sum of the totient]\label{el:phisum}
For every $n\ge1$,
\[
                         \sum_{d\mid n}\varphi(d)=n.
\]
\end{theorem}
\begin{proof}
Partition $a=1,\ldots,n$ by $g=\gcd(a,n)$.  For a fixed divisor $g$
of $n$, the assignment $a=gu$ identifies this part with the integers
$1\le u\le n/g$ coprime to $n/g$: dividing a B\'ezout identity for
$a,n$ by $g$ proves one direction, and multiplying a coprimality
identity for $u,n/g$ by $g$ proves the other.  This part has
$\varphi(n/g)$ members.  The divisor map $g\mapsto n/g$ is an
involution, so counting all $n$ integers gives the stated sum.
This is a disjoint finite classification, implementable by a double
scan bounded by $n^2$.
\end{proof}

\begin{theorem}[M\"obius cancellation]\label{el:mobiuscancel}
The function $\mu$ is multiplicative and
\[
 \sum_{d\mid n}\mu(d)=\varepsilon(n),\qquad
 \varepsilon(n)=\begin{cases}1&n=1,\\0&n>1.\end{cases}
\]
\end{theorem}
\begin{proof}
For coprime inputs, their prime supports are disjoint.  If either has
a repeated prime factor, the product does too and both sides of
multiplicativity are zero.  Otherwise the numbers of prime factors
add and their parity signs multiply.  This proves multiplicativity.
For $n=1$ the divisor sum is $\mu(1)=1$.  For $n>1$, fix the least
prime $p\mid n$.  Divisors with a squared prime factor contribute
zero.  Pair each squarefree divisor not containing $p$ with its
product by $p$; both divide $n$, the pairing is reversible, and the
values of $\mu$ are opposite.  The finite sum is therefore zero.
Equivalently, distributing over the distinct prime support gives a
product of factors $1-1$.  The pairing is explicit and bounded by $n$.
\end{proof}

\begin{theorem}[Liouville's function and square detection]
\label{el:liouville}
Let $\Omega(n)$ be the length of the prime factor list, counting
multiplicity, and put $\lambda(n)=(-1)^{\Omega(n)}$ for $n\ge1$.
Then $\lambda(ab)=\lambda(a)\lambda(b)$ for all positive $a,b$, and
\[
 \sum_{d\mid n}\lambda(d)=
 \begin{cases}1&n\text{ is a square},\\0&n\text{ is not a square}.
 \end{cases}
\]
\end{theorem}
\begin{proof}
Concatenating prime factor lists and using their uniqueness gives
$\Omega(ab)=\Omega(a)+\Omega(b)$, without any coprimality assumption.
This proves complete multiplicativity of $\lambda$.  A divisor of
$n=\prod_p p^{e_p}$ chooses each exponent independently, so finite
distributivity gives
\[
 \sum_{d\mid n}\lambda(d)=\prod_{p\mid n}
             (1-1+1-\cdots+(-1)^{e_p}).
\]
Each parenthesis is one when $e_p$ is even and zero when it is odd,
by pairing consecutive terms.  The product is therefore one exactly
when every exponent is even.  This is the square criterion just proved.
For $n=1$ the empty product is one, as required.
\end{proof}

For fixed integer-valued primitive recursive functions $f,g$ define
Dirichlet convolution by
\[
                  (f*g)(n)=\sum_{d\mid n} f(d)g(n/d),\qquad n\ge1.
\]
This is a primitive recursive operation for each fixed pair of function
definitions.  One can instead take arrays containing $f(1),\ldots,f(N)$
and $g(1),\ldots,g(N)$ to obtain a single first-order assertion for
all indices at most $N$.

\begin{theorem}[Finite Dirichlet convolution algebra]
\label{el:convolution}
Dirichlet convolution is commutative and associative, has identity
$\varepsilon$, and distributes over pointwise addition.  If $f,g$
are multiplicative, then $f*g$ is multiplicative.
\end{theorem}
\begin{proof}
For commutativity replace a divisor $d$ by $n/d$; this substitution
is a bijection of the finite divisor list.  Both
$((f*g)*h)(n)$ and $(f*(g*h))(n)$ expand as
\[
                   \sum_{abc=n}f(a)g(b)h(c),
\]
where $1\le a,b,c\le n$.  The equality follows by a bijective
reindexing of this bounded triple sum, proving associativity.
Only the divisor one contributes to $(\varepsilon*f)(n)$, and
finite distributivity proves the additive rule.

If $\gcd(a,b)=1$, each divisor of $ab$ has a unique form $uv$ with
$u\mid a$ and $v\mid b$, by the prime exponent description.  The
complementary divisor is $(a/u)(b/v)$.  The two factors in each of
these products are coprime.  Multiplicativity of $f$ and $g$ therefore
gives
\[
 (f*g)(ab)=\sum_{u\mid a}\sum_{v\mid b}
 f(u)f(v)g(a/u)g(b/v)=(f*g)(a)(f*g)(b).
\]
The value at one is one.  For a formal theorem about finite tables,
only the multiplicativity identities at the bounded arguments used
in this display are hypotheses; no quantification over functions is
part of the arithmetic language.
\end{proof}

\begin{theorem}[M\"obius inversion]\label{el:mobius}
For fixed primitive recursive integer-valued $f,F$, the identities
\[
 F(n)=\sum_{d\mid n}f(d)\quad(n\ge1)
\]
are equivalent to
\[
 f(n)=\sum_{d\mid n}\mu(d)F(n/d)\quad(n\ge1).
\]
For each $N$, the same equivalence holds using only the corresponding
identities at indices $1,\ldots,N$.  In particular
\[
                   \varphi(n)=\sum_{d\mid n}\mu(d)(n/d).
\]
\end{theorem}
\begin{proof}
Let $\mathbf1(n)=1$.  M\"obius cancellation says
$\mu*\mathbf1=\varepsilon$.  If $F=\mathbf1*f$, associativity gives
$\mu*F=(\mu*\mathbf1)*f=f$.  Conversely, if $f=\mu*F$, then
$\mathbf1*f=(\mathbf1*\mu)*F=F$.
Every convolution evaluated at $n$ uses only divisors of $n$, all at
most $n$, so the proof applies verbatim to the bounded version with
$n\le N$.  Apply it to Theorem~\ref{el:phisum} to obtain the totient
formula.  This proof is therefore a quantifier-free induction on
finite sums after the bounded hypotheses have been folded into one
Boolean antecedent, rather than an induction on a functional equation.
\end{proof}

\begin{theorem}[Divisor summatory and hyperbola identities]
\label{el:summatory}
For $N\in\N$ and $r\in\N$,
\[
 \sum_{n=1}^N\sigma_r(n)=
      \sum_{d=1}^N d^r\left\lfloor\frac Nd\right\rfloor.
\]
Writing $h=\lfloor\sqrt N\rfloor$, the divisor-counting case also has
\[
 \sum_{n=1}^N\tau(n)=
       2\sum_{d=1}^{h}\left\lfloor\frac Nd\right\rfloor-h^2.
\]
\end{theorem}
\begin{proof}
Both sides of the first identity count the finite pairs $(d,k)$ with
$d,k\ge1$ and $dk\le N$, giving pair weight $d^r$.  On the left
classify by $n=dk$; on the right fix $d$ and count
$k\le\lfloor N/d\rfloor$.  This is a reindexing in the square
$[1,N]^2$.  In the unweighted case every admissible pair has $d\le h$
or $k\le h$, because otherwise $dk\ge(h+1)^2>N$.  The two regions
each have $\sum_{d=1}^{h}\lfloor N/d\rfloor$ points.  Their intersection
is the full square $[1,h]^2$, since $h^2\le N$, and has $h^2$ points.
Adding the two region counts and subtracting their intersection gives
the second formula.  For $N=0$ all regions and sums are empty.
\end{proof}

\begin{theorem}[Euclid--Euler characterization of even perfect numbers]
\label{el:perfect}
An even positive integer $N$ satisfies $\sigma(N)=2N$ exactly when
\[
                  N=2^{e-1}(2^e-1)
\]
for some $e\ge2$ such that $2^e-1$ is prime.  The exponent $e$ is then
prime.  For a supplied $N$ all witnesses and all tests are bounded
primitive recursive computations.
\end{theorem}
\begin{proof}
If $q=2^e-1$ is prime and $e\ge2$, then $q$ is odd, and
multiplicativity of $\sigma$ gives
\[
 \sigma(2^{e-1}q)=(1+2+\cdots+2^{e-1})(1+q)
                  =(2^e-1)2^e=2N.
\]
Conversely, write an even $N$ as $2^s m$ with $s\ge1$ and $m$ odd,
using $s=v_2(N)$.  The equation for perfection becomes
\[
             (2^{s+1}-1)\sigma(m)=2^{s+1}m.
\]
Put $q=2^{s+1}-1$.  Since $q$ and $2^{s+1}$ are consecutive, they are
coprime, so coprime cancellation gives $q\mid m$.  Write $m=qt$;
then the displayed equation simplifies to
$\sigma(m)=(q+1)t=m+t$.  If $t>1$, the three distinct numbers
$1,t,m$ are divisors of $m$ (here $q\ge3$), so
$\sigma(m)\ge1+t+m>t+m$, a contradiction.  Thus $t=1$ and $m=q$.
The identity $\sigma(q)=q+1$ forces $q$ to be prime, because a proper
divisor $1<d<q$ would add an additional positive summand besides
$1$ and $q$.  Taking $e=s+1$ gives the required form.
Finally, if $e=uv$ with $u,v>1$, then
\[
 2^e-1=(2^u-1)(1+2^u+\cdots+2^{u(v-1)})
\]
has two factors greater than one, contradicting primality.
The exponent $s$ is obtained by a scan bounded by $N$, and the proof
uses only factorization, finite sums, and explicit inequalities.  It
says nothing about the existence or nonexistence of odd perfect numbers.
\end{proof}

\subsection{Quantifier-free proof obligations for this chapter}

The constructions above provide witness terms, not merely assertions that
an algorithm eventually terminates.  Their common verification pattern is
as follows.  Code a loop state by a finite tuple of natural numbers; for
signed coordinates, code a sign and magnitude.  Set a total clock, extend
the transition by the identity after its stopping condition, and define
the state after $i$ steps by primitive recursion.  Express the invariant
as equality of natural-valued characteristic functions.  Its preservation
is a quantifier-free statement with the original inputs as parameters.
Apply quantifier-free induction to $i$, then decode the state at the clock.
Finite permutations and reindexings are verified by explicit bounded
maps and their inverses.  No infinite set, quotient type, arbitrary
choice function, or unbounded minimization is needed.

For clarity, some useful clocks and magnitude bounds are collected here.
They are deliberately generous; optimal bit complexity is a separate
question.
\begin{center}
\begin{tabular}{p{0.27\linewidth}p{0.62\linewidth}}
\hline
Construction & Sufficient explicit arithmetic bound\\
\hline
Division of $a$ by $b>0$ & $a$ counter steps; outputs at most $a$.\\
Extended Euclid & $a+b+1$ steps; coefficient magnitudes at most
$(a+b+2)^{a+b+1}$.\\
Prime factorization & $n$ division stages, each divisor scan at most $n$;
factor array of length at most $n$ with entries at most $n$.\\
Sieve to $N$ & At most $N^2$ crossing-off checks; bit-array code
below $2^{N+1}$.\\
CRT for $k$ moduli & Representative below $\prod_{i<k}m_i$;
prime and exponent scans below that product.\\
Orders and inverses & Residue scan below $n$ and exponent scan
at most $\varphi(n)\le n$.\\
Arithmetic functions & Divisor and prime scans bounded by $n$;
$\sigma_r(n)\le n^{r+1}$ for $n\ge1$.\\
Summatory identities & Reindexings inside $[1,N]^2$.\\
\hline
\end{tabular}
\end{center}

Each theorem concerning a positive input is encoded with a decidable
antecedent, so total definitions may return a harmless default on
nonpositive or malformed inputs.  This handles zero denominators and
invalid finite-array codes without adding a partial function to \ipra.
